\section{Guion de entrevista estructurada}
\label{sec:guion-entrevista}

Se presenta un guion de trece preguntas para un ensayo de entrevista estructurada con una persona que representa el rol de cajero de un establecimiento de productos eléctricos de Quevedo. El objetivo es conocer el proceso de venta, la facturación, los posibles errores de registro y las consultas de información técnica antes de proponer funciones para la plataforma de ventas. La secuencia avanza desde las actividades generales hasta situaciones específicas, mediante preguntas principalmente abiertas y neutrales, conforme a las recomendaciones metodológicas de elicitación de requisitos \cite{pohl2025,ireb2026}. Las respuestas del juego de roles no constituyen evidencia de entrevistas con trabajadores reales.

\subsection{Apertura}
Antes de las preguntas se explica el propósito académico, la duración aproximada de diez minutos y la participación voluntaria. El entrevistador solicita consentimiento explícito para grabar únicamente el audio. Si no se autoriza la grabación, esta no se realiza.

\noindent\textbf{Presentación sugerida:} ``Buenos días. Esta actividad académica busca comprender cómo se desarrolla la atención y la facturación en un establecimiento de productos eléctricos. No evaluaremos su desempeño personal. La entrevista durará aproximadamente diez minutos; puede omitir una pregunta o detenerse en cualquier momento. Solo grabaremos audio si usted lo autoriza.''

\subsection{Proceso de venta}
Las preguntas 3 y 4 permiten describir la secuencia de atención al cliente y la forma de localizar los productos solicitados, sin asumir el uso de un sistema específico.

\subsection{Facturación}
Las preguntas 5 y 6 indagan el registro de los productos y la generación de facturas, con el propósito de reconocer pasos y herramientas actuales.

\subsection{Dificultades operativas}
La pregunta 7 identifica pasos que podrían requerir más tiempo, dejando abierta la posibilidad de que el entrevistado no observe demoras significativas.

\subsection{Errores de digitación}
Las preguntas 8 y 9 investigan experiencias de errores de registro y sus procedimientos de corrección, sin presumir que estos hayan ocurrido.

\subsection{Atención en momentos de alta demanda}
La pregunta 10 busca describir cambios en la actividad del cajero cuando varios clientes esperan atención, sin atribuir previamente dificultades.

\subsection{Consulta de información técnica}
La pregunta 11 identifica cómo se obtiene información sobre las características de los productos eléctricos cuando se necesita consultarla.

\subsection{Orientación sobre productos}
La pregunta 12 examina cómo se atienden las solicitudes de orientación de los clientes sin asumir que el cajero realiza recomendaciones técnicas por cuenta propia.

\subsection{Cierre}
La pregunta 13 permite incorporar aspectos que no se hayan tratado. Se agradece la participación y se finaliza la grabación, si fue autorizada.

\noindent\textbf{Despedida sugerida:} ``Gracias por compartir su experiencia. Sus respuestas nos ayudarán a comprender mejor las actividades de venta, facturación y consulta de productos. Agradecemos su participación en esta práctica.''

\begingroup
\small
\renewcommand{\arraystretch}{1.17}
\setlength{\LTpre}{9pt}
\setlength{\LTpost}{8pt}
\begin{longtable}{@{}p{0.045\textwidth}@{\hspace{7pt}}p{0.139\textwidth}@{\hspace{7pt}}p{0.385\textwidth}@{\hspace{7pt}}p{0.105\textwidth}@{\hspace{7pt}}p{0.235\textwidth}@{}}
\caption{Preguntas del guion de entrevista estructurada al cajero}\label{tab:guion-entrevista}\\
\toprule
\textbf{N.º} & \textbf{Bloque} & \textbf{Pregunta} & \textbf{Tipo} & \textbf{Dato que busca obtener}\\
\midrule
\endfirsthead
\multicolumn{5}{@{}l}{\textit{Tabla \thetable\ (continuación)}}\\
\toprule
\textbf{N.º} & \textbf{Bloque} & \textbf{Pregunta} & \textbf{Tipo} & \textbf{Dato que busca obtener}\\
\midrule
\endhead
\midrule
\multicolumn{5}{r}{\textit{Continúa en la siguiente página}}\\
\endfoot
\bottomrule
\endlastfoot
1 & Apertura & ¿Autoriza que grabemos únicamente el audio de esta entrevista académica? & Cerrada & Consentimiento para la grabación.\\
2 & Apertura & ¿Cuáles son sus principales actividades como cajero durante una jornada de trabajo? & Abierta & Responsabilidades y funciones habituales.\\
3 & Proceso de venta & ¿Cómo se desarrolla normalmente una venta desde que comienza la atención al cliente? & Abierta & Secuencia general del proceso de venta.\\
4 & Proceso de venta & ¿Cómo localiza el producto que solicita un cliente? & Abierta & Procedimiento de búsqueda de productos.\\
5 & Facturación & ¿Cómo registra los productos que forman parte de una venta? & Abierta & Procedimiento de ingreso de productos.\\
6 & Facturación & ¿Cómo genera actualmente la factura de una venta? & Abierta & Pasos y herramientas de facturación.\\
7 & Dificultades operativas & ¿Qué pasos de la facturación le toman más tiempo, si los hay? & Abierta & Posibles actividades que generan demoras.\\
8 & Errores de digitación & ¿Qué experiencias ha tenido, si alguna, con errores al ingresar datos de una venta? & Abierta & Errores experimentados u observados.\\
9 & Errores de digitación & Si ha ocurrido un error al registrar una venta, ¿cómo lo corrigieron? & Abierta & Procedimiento real de corrección.\\
10 & Alta demanda & ¿Cómo cambia su trabajo cuando hay varios clientes esperando atención? & Abierta & Efectos de la demanda sobre las actividades.\\
11 & Información técnica & ¿Cómo obtiene la información técnica de un producto cuando necesita consultarla? & Abierta & Fuentes de información utilizadas.\\
12 & Recomendaciones & ¿Qué hace cuando un cliente le pide orientación para elegir un producto eléctrico? & Abierta & Proceso de atención de consultas técnicas.\\
13 & Cierre & ¿Qué otro aspecto de su trabajo considera importante mencionar? & Abierta & Necesidades adicionales no identificadas.\\
\end{longtable}
\endgroup

\noindent\textbf{Distribución prevista del tiempo:} apertura y consentimiento (1 min 30 s); proceso de venta (1 min 20 s); facturación (1 min 25 s); dificultades y errores (3 min); información técnica y orientación (1 min 40 s); cierre (50 s), con 15 s para transiciones. Duración total prevista: diez minutos.

\noindent\textbf{Sondeos neutrales:} cuando se necesite ampliar una respuesta, se podrá preguntar ``¿Podría darme un ejemplo?'' o ``¿Qué ocurrió después?'' sin sugerir la respuesta. Si el tiempo se reduce a siete minutos, se priorizan las preguntas 1, 2, 3, 5, 7, 8, 10, 11, 12 y 13.
